\documentclass[12pt,a4paper,UTF8]{ctexart}
\usepackage[a4paper,top=22mm,bottom=22mm,left=22mm,right=22mm,headheight=15pt]{geometry}
\newcommand{\Needspace}[1]{\par\begingroup\dimen0=#1\relax\dimen2=\pagegoal\advance\dimen2 by-\pagetotal\ifdim\dimen0>\dimen2\newpage\fi\endgroup}
\newcommand{\needspace}[1]{\par\begingroup\skip0=\lastskip\vskip-\skip0\dimen0=#1\relax\advance\dimen0 by\skip0\vskip\dimen0\penalty-100\vskip-\dimen0\vskip\skip0\penalty9999\endgroup}
\usepackage{amsmath,amssymb,mathtools,bm}
\usepackage{graphicx,xcolor,booktabs,longtable,array,enumitem,fancyhdr,caption,float,microtype,setspace}
\usepackage[hidelinks,unicode]{hyperref}
\usepackage[most]{tcolorbox}
\definecolor{ContestBlue}{HTML}{17365D}
\hypersetup{pdftitle={第059套},pdfauthor={}}
\setstretch{1.24}
\setlength{\parindent}{2em}
\setlength{\parskip}{0.35em}
\setlength{\emergencystretch}{2em}
\setlist{itemsep=0.35em,topsep=0.4em}
\captionsetup{font=small,labelfont=bf,justification=centering}
\setcounter{secnumdepth}{0}
\pagestyle{fancy}
\fancyhf{}
\fancyhead[L]{\small 第059套}
\fancyfoot[C]{\thepage}
\renewcommand{\headrulewidth}{0.4pt}
\ctexset{
  section={format=\Large\bfseries\color{ContestBlue},beforeskip=1.2em,afterskip=0.7em},
  subsection={format=\large\bfseries\color{ContestBlue},beforeskip=1.0em,afterskip=0.5em},
  subsubsection={format=\normalsize\bfseries,beforeskip=0.8em,afterskip=0.35em}
}
\usepackage{calc}
\newcounter{none}
\let\originalsubsection\subsection
\renewcommand{\subsection}{\needspace{6\baselineskip}\originalsubsection}
\makeatletter
\newcommand{\tightlist}{\setlength{\itemsep}{0pt}\setlength{\parskip}{0pt}}
\newcommand{\maxwidth}{\ifdim\Gin@nat@width>\linewidth\linewidth\else\Gin@nat@width\fi}
\newcommand{\maxheight}{\ifdim\Gin@nat@height>0.72\textheight 0.72\textheight\else\Gin@nat@height\fi}
\newcommand{\pandocbounded}[1]{\begingroup\setkeys{Gin}{width=\maxwidth,height=\maxheight,keepaspectratio}#1\endgroup}
\makeatother
\begin{document}
\begin{titlepage}
  \centering
  \vspace*{58mm}
  {\Huge\bfseries\color{ContestBlue} 第059套\par}
  \vfill
  \vspace*{28mm}
\end{titlepage}
\section{第 1 题}\label{ux7b2c-1-ux9898}

海湾呈顶点为 \(O\)、顶角为锐角 \(\alpha\) 的 V 形，角 \(\angle AOB\)
内为海水，外部为陆地。目标位于海岸 \(OA\) 上的点 \(D\)，且
\(OD=x\)。起点为 \(H\)，其在海岸 \(OB\) 上的垂足为
\(C\)，\(OC=l\)、\(CH=h\)。人在陆地上的奔跑速度为
\(v\)，在静水中的游泳速度为 \(v/n\)。

已知

\[
n=2,\qquad l=x=1\,\mathrm{km},\qquad
h=0.5\,\mathrm{km},\qquad
\alpha=\frac{\pi}{4},\qquad v=1\,\mathrm{m/s}.
\]

（1）若全程不下水，求最短用时。

（2）从 \(H\) 直线跑到 \(OB\) 上一点 \(E\)，再从 \(E\) 直线游到
\(D\)。求最短用时及 \(OE\)。

（3）从 \(H\) 直线跑到 \(OB\) 上一点 \(F\)，从 \(F\) 直线游到 \(OA\)
上一点 \(G\)，再沿岸跑到 \(D\)。求最短用时及 \(OF,OG\)。

（4）比较三种方案，指出最快者。

\begin{figure}[H]
\centering
\includegraphics[width=0.564\linewidth,height=\textheight,keepaspectratio,alt={图1-1　V 形海湾、起点与目标点的几何位置}]{../figures/question-01.pdf}
\caption*{图1-1　V 形海湾、起点与目标点的几何位置}
\end{figure}

\newpage

\section{第 2 题}\label{ux7b2c-2-ux9898}

两个质量均为 \(m\) 的均匀弹性小球，初始质心处于同一高度，水平间距为
\(l\)，由一根长为 \(3l\) 的不可伸长轻绳连接，装置位于距地面足够高处。

（1）若绳均匀、总质量为
\(M\)，在重力作用下自由下垂。绳与小球的连接点比绳最低点高
\(h\)。求绳上的最大张力。

（2）忽略绳质量，解除维持初始位置的支承，让绳保持松弛且仍连接两球。左球自由下落，同时给右球一个指向左球、大小为
\(v_0\) 的水平初速度。求从开始运动到两球相碰期间左球下降的高度。

（3）两球发生完全弹性碰撞，求碰撞后各球速度的水平分量。

（4）碰后绳逐渐被拉直。求轻绳拉直的短暂过程中，右球受到绳张力的冲量大小。

\begin{figure}[H]
\centering
\includegraphics[width=0.78\linewidth,height=\textheight,keepaspectratio,alt={图2-1　等高小球与自由下垂均匀绳的几何关系}]{../figures/question-02.pdf}
\caption*{图2-1　等高小球与自由下垂均匀绳的几何关系}
\end{figure}

\newpage

\section{第 3 题}\label{ux7b2c-3-ux9898}

平面内有一根按对数螺线弯成的均匀细铁丝，其极坐标方程为

\[
\rho=r_0e^{k\theta},\qquad
k>0,\qquad \theta\in(-\infty,0).
\]

铁丝线密度为 \(\lambda\)。

（1）求铁丝总质量。

（2）利用``在末端增加微元后，图形可由缩放与微小旋转得到''的自相似方法，求质心坐标
\((x_c,y_c)\)。

（3）用同样的自相似思想求铁丝关于原点的转动惯量。

（4）取
\(k=1\)。铁丝可绕原点无摩擦转动，初始静止。一只初始也静止、质量与铁丝总质量相等的虫沿铁丝缓慢爬向原点。求虫到达原点的极限过程中铁丝转过的角度。

\begin{figure}[H]
\centering
\includegraphics[width=0.78\linewidth,height=\textheight,keepaspectratio,alt={图3-1　对数螺线铁丝的极坐标定义与末端半径}]{../figures/question-03.pdf}
\caption*{图3-1　对数螺线铁丝的极坐标定义与末端半径}
\end{figure}

\newpage

\section{第 4 题}\label{ux7b2c-4-ux9898}

点电荷 \(Q\) 位于 \(z\) 轴上、距原点为 \(L\)
的位置。原点附近放置一根长为 \(2r\) 的细金属棒，且
\(2r\ll L\)。棒初始中性，系统达到静电平衡。前两问先按忽略横向极化的细线近似求幂次；有限粗细的修正在解析中另行说明。

（1）金属棒沿 \(z\) 轴放置。若点电荷受到的作用力写成

\[
F=A_1Q^{\alpha_1}L^{\beta_1}r^{\gamma_1},
\]

求 \(\alpha_1,\beta_1,\gamma_1\)。

（2）金属棒改沿 \(x\) 轴放置。若

\[
F=A_2Q^{\alpha_2}L^{\beta_2}r^{\gamma_2},
\]

求 \(\alpha_2,\beta_2,\gamma_2\)。

（3）已知焦距为 \(c\)、离心率为 \(e\)、均匀体电荷密度为 \(\rho\)
的长轴旋转椭球，其内部电势可写成

\[
\phi(x,y,z)
=\rho\frac{b^2}{c^2}
\left[\Gamma_1(e)(x^2+y^2)+\Gamma_2(e)z^2\right],
\]

其中 \(\Gamma_1,\Gamma_2\)
为已知函数。利用细长椭球近似金属棒，求第（1）问中的系数 \(A_1\)。

\begin{figure}[H]
\centering
\includegraphics[width=0.855\linewidth,height=\textheight,keepaspectratio,alt={图4-1　点电荷与纵向、横向细金属棒的两种构形}]{../figures/question-04.pdf}
\caption*{图4-1　点电荷与纵向、横向细金属棒的两种构形}
\end{figure}

\newpage

\section{第 5 题}\label{ux7b2c-5-ux9898}

质量为 \(m\)、半径为 \(r\) 的均匀实心球在水平面上运动，转动惯量

\[
I=\frac25mr^2.
\]

取竖直向上为 \(+z\)，水平面内建立 \(xOy\) 坐标。以下 \(z\)
表示球心相对初始高度的位移；实际地面位于球心下方 \(r\)。初始条件为

\[
x(0)=y(0)=z(0)=0,
\]

\[
v_x(0)=v_0=r\omega_0,\qquad
v_y(0)=v_z(0)=0,
\]

\[
\omega_x(0)=0,\qquad
\omega_y(0)=\omega_0,\qquad
\omega_z(0)=0.
\]

空间中存在 \(\bm B=-B\hat{\bm z}\) 的匀强磁场。

（1）电荷 \(q\) 固定在球心，地面静摩擦足够大，求球心的运动。

（2）求维持上述纯滚动所需的临界静摩擦因数
\(\mu_{\max}\)；当地面摩擦因数低于它时可能滑动。

（3）若总电荷 \(q\) 均匀分布在球面，球以角速度 \(\omega_0\)
转动，求球的磁矩大小。

（4）承接第（3）问，地面与球均绝缘且静摩擦足够大，求球心的运动。

\begin{figure}[H]
\centering
\includegraphics[width=0.82\linewidth,height=\textheight,keepaspectratio,alt={图5-1　带电均匀球的初始纯滚动状态与磁场方向}]{../figures/question-05.pdf}
\caption*{图5-1　带电均匀球的初始纯滚动状态与磁场方向}
\end{figure}

\newpage

\section{第 6 题}\label{ux7b2c-6-ux9898}

角频率为 \(\omega\) 的平面电磁波在真空中沿 \(+x\)
方向传播。旋转抛物反射面的瞬时方程为

\[
x=l-\alpha(y^2+z^2)-v_0t,
\]

它以速度 \(v_0\) 沿 \(-x\) 方向运动。

（1）在实验室系中，求时刻 \(t\) 的聚焦点横坐标 \(x_1(t)\)。

（2）求实验室系用坐标时间描述的聚焦点处角频率 \(\omega_1\)。

\begin{figure}[H]
\centering
\includegraphics[width=0.9\linewidth,height=\textheight,keepaspectratio,alt={图6-1　迎面运动的旋转抛物反射面与入射平面波}]{../figures/question-06.pdf}
\caption*{图6-1　迎面运动的旋转抛物反射面与入射平面波}
\end{figure}

\newpage

\section{第 7 题}\label{ux7b2c-7-ux9898}

截面积为 \(a\)、总长为 \(2L\)
的绝热气缸，被一个导热性很差、可无摩擦移动的光滑活塞分成左右两部分。初始时左右体积之比为
\(1:3\)，两侧各装有等量的双原子理想气体。左侧初压为
\(3^\gamma p_0\)，右侧初压、初温分别为 \(p_0,T_0\)，其中

\[
\gamma=\frac{C_p}{C_v}=\frac75.
\]

以气缸中点为坐标原点，沿气缸向右为 \(x\)
轴正方向。导热速率和活塞运动很缓慢，使气体始终处于准静态；活塞单次振动周期
\(t_2\) 远小于左右气体达到热平衡的时间 \(t_3\)，而气体内部弛豫时间
\(t_1\) 远小于 \(t_2\)。

（1）从初始位置由静止释放活塞。活塞第一次速度为零时的位置为
\(x_1\)，定义 \(x_1/L=\alpha_1\)，求 \(\alpha_1\)。

（2）求最终热平衡时活塞位置 \(x_\infty/L=\alpha_\infty\)。

（3）求最终平衡温度 \(T_\infty/T_0=\beta\)。

（4）若系统总熵增为 \(\Delta S\)，求

\[
\eta=\frac{\Delta S\,T_0}{p_0aL}.
\]

（5）把活塞相邻两次速度为零的位置间距定义为振幅。初始振幅为
\(A_0\)；当振幅衰减到
\(A_{0.5}=A_0/2\)，且活塞位于左侧转折点时，左侧气体温度满足
\(T_{0.5}/T_0=\xi\)。求 \(\xi\)。

\begin{figure}[H]
\centering
\includegraphics[width=0.82\linewidth,height=\textheight,keepaspectratio,alt={图7-1　绝热气缸、低导热活塞与两侧气体的初始状态}]{../figures/question-07.pdf}
\caption*{图7-1　绝热气缸、低导热活塞与两侧气体的初始状态}
\end{figure}

\newpage

\section{第 8 题}\label{ux7b2c-8-ux9898}

在呼啦圈的静止参考系中，呼啦圈半径为 \(R\)。

（1）以呼啦圈为静止参考系 \(S\)。参考系 \(S'\) 以速度 \(v\) 沿 \(-x\)
方向运动。在 \(S'\) 中同时测量呼啦圈，求所得图形的面积 \(A_1\)。

（2）呼啦圈绕中心沿图示顺时针方向以角速度 \(\omega\) 转动，且
\(\omega R<c\)。圈上每一点都在其自身瞬时静止系中向所有方向发出频率为
\(\nu\) 的光。在 \(S\) 系中，位于 \((0,y)\) 的接收器接收到由角位置
\(\theta_1\) 处发出的光，求接收频率 \(\nu_1\)。

（3）沿用第（2）问的装置。两参考系原点重合时，各自把时钟调为
\(t=t'=0\)。在 \(S'\) 系中的 \((0,y')\) 处，接收器于 \(t'=0\)
收到从圈上角位置 \(\theta\) 处的点 \(P\) 发出的光。求该光发出事件在
\(S'\) 中的坐标 \((x'_P,y'_P)\)。

\begin{figure}[H]
\centering
\includegraphics[width=0.74\linewidth,height=\textheight,keepaspectratio,alt={图8-1　呼啦圈、转动方向及两个参考系的坐标约定}]{../figures/question-08.pdf}
\caption*{图8-1　呼啦圈、转动方向及两个参考系的坐标约定}
\end{figure}
\end{document}
